\documentclass[sigconf]{acmart}
\setlength{\emergencystretch}{1.5em}

% AISec 2026 follows the ACM CCS 2026 formatting instructions, except for
% its 10-page body + 2-page appendix limit.  The CCS 2026 minimal sample
% requires these provisional rights-management fields at submission time.
\setcopyright{acmlicensed}
\copyrightyear{2018}
\acmYear{2018}
\acmDOI{XXXXXXX.XXXXXXX}
\acmConference[AISec '26]{Proceedings of the 2026 ACM Workshop on Artificial
  Intelligence and Security}{November 15--19, 2026}{The Hague, Netherlands}
\acmISBN{978-1-4503-XXXX-X/2018/06}

\title{Detection Is Not Defense: Stress-Testing Spectral Monitoring for Hidden-State Poisoning}

\begin{CCSXML}
<ccs2012>
<concept>
<concept_id>10002978.10002997</concept_id>
<concept_desc>Security and privacy~Intrusion/anomaly detection and malware mitigation</concept_desc>
<concept_significance>500</concept_significance>
</concept>
<concept>
<concept_id>10010147.10010257</concept_id>
<concept_desc>Computing methodologies~Machine learning</concept_desc>
<concept_significance>300</concept_significance>
</concept>
</ccs2012>
\end{CCSXML}

\ccsdesc[500]{Security and privacy~Intrusion/anomaly detection and malware mitigation}
\ccsdesc[300]{Computing methodologies~Machine learning}
\keywords{state space models, machine learning security, adversarial machine learning, model robustness}
\begin{document}
\begin{abstract}
Hidden-state poisoning can corrupt the recurrent memory of state-space
language models, but a detector-level score does not establish preserved task
behavior or operational availability. We first repair the theoretical basis
for spectral monitoring: for frozen diagonal transitions we derive an exact
forgetting certificate, and a constructive study of 1,000 diagonal systems
shows that readout failure can coexist with an unchanged spectral radius,
ruling out unconditional completeness for radius-only monitoring. We then
stress-test layer-wise spectral detection under a leakage-controlled protocol
spanning seven models, four tasks, and 19,300 evaluated rows. SpectralGuard
obtains F1 0.806 and AUC 0.887, whereas a CLASP-style BOE+XGBoost comparator
obtains mean F1 0.969 and AUC 0.997 under the same internal split protocol.
Detector-aware adaptive attacks reduce attacked task accuracy to 1.67\%,
statistically indistinguishable from 1.33\% for the preceding adaptive attack
in a floor-saturated comparison; standardized mean spectral radius alone has
weak explanatory power for correctness (McFadden pseudo-$R^2$
0.0089--0.0337). Finally, a fail-closed policy blocks all 320 held-out clean
sessions, while a validation-selected graduated policy allows only 8/320 to
survive and incurs 4.29$\times$ total latency. These results show that spectral
features carry predictive signal, but neither high aggregate detection scores
nor radius-only mechanisms establish a robust and available defense.
Hidden-state-poisoning evaluations must jointly test detection, attacked-task
utility, adaptive evasion, and session-level availability.
\end{abstract}
\maketitle

% Double-blind submission: deliberately no \author, \affiliation, \email,
% acknowledgments, funding statement, or identifying repository URL appears.

\section{Introduction}
State-space language models (SSMs) offer linear-time sequence processing by
maintaining a recurrent hidden state. That state is also a security boundary:
a short hostile insertion can alter later behavior without remaining in the
visible context. Hidden State Poisoning Attacks (HiSPAs) demonstrate this
failure mode in Mamba-family models \cite{lemercier2026hispa}. A natural
response is to monitor internal dynamics, but a high detector score answers
only whether selected attacks are separable on a benchmark. It does not show
that useful behavior survives an adaptive attack, that a scalar statistic is a
complete failure indicator, or that false alarms permit benign sessions to
continue.

We examine that gap through SpectralGuard, a layer-wise monitor derived from
the discretized state transition. The study began with a stronger hypothesis:
that spectral contraction explained hidden-state poisoning and enabled a
complete detector. Our audit rejects that hypothesis. Spectral radius gives an
exact forgetting upper bound for frozen diagonal transitions, yet an unchanged
radius can coexist with total readout loss. Empirically, radius is predictive
but weakly explanatory, and a richer spectral detector is substantially
outperformed by a CLASP-style block-output-embedding comparator under one
shared internal protocol. Most importantly, adaptive attacks preserve very
low attacked-task accuracy, while a fail-closed deployment rule denies service
to every clean held-out session.

This paper makes five contributions:
\begin{itemize}
  \item We replace an invalid persistence argument with an exact diagonal
  forgetting certificate and state its assumptions explicitly.
  \item We provide 1,000 constructive counterexamples showing that a
  radius-only monitor is not unconditionally complete.
  \item We evaluate SpectralGuard across seven models, four tasks, and 19,300
  rows with disjoint calibration, validation, and test roles.
  \item We compare against a CLASP-style BOE+XGBoost implementation and
  distinguish fixed-test-set bootstrap uncertainty from seed-to-seed training
  variation; we make no detector-superiority claim.
  \item We connect detection to attacked-task utility, adaptive evasion, and
  multi-turn availability, revealing a defense-induced denial-of-service
  failure that aggregate F1 and AUC conceal.
\end{itemize}

\section{Background and Related Work}

Le Mercier, Develder, and Demeester introduced Hidden State Poisoning
Attacks (HiSPAs), showing that short adversarial strings can corrupt the
recurrent memory of Mamba-family models and evaluating the effect with the
RoBench-25 benchmark \cite{lemercier2026hispa}.  Their latest arXiv revision
also studies Mamba-2 and hybrid SSM--Transformer models.

Concurrent work proposes CLASP, a token- and document-level detector trained
with XGBoost over Mamba block-output embeddings (BOEs)
\cite{lemercier2026clasp}.  CLASP is evaluated primarily on controlled
injections into a resume corpus and explicitly leaves adaptive evasion and
cross-domain validation open.  SpectralGuard instead investigates the
dynamics of the discretized state-transition operator and a detector based on
layer-wise spectral measurements.  Under the common Phase~3 protocol,
SpectralGuard obtained F1~$0.806$ and AUC~$0.887$, while our CLASP-style
BOE+XGBoost reimplementation obtained F1~$0.969$ and AUC~$0.997$.  The latter
uses the published CLASP design elements rather than official pretrained
weights; accordingly, we claim neither that the reimplementation is identical
to the authors' release nor that SpectralGuard is empirically superior.

Adaptive evaluation must also test whether a defense relies on brittle or
obfuscated gradients \cite{athalye2018obfuscated}. More broadly, security
systems can impose operational costs that benchmark-level false-positive rates
hide. Alahmadi et al. quantify the cost of false positives in network intrusion
detection \cite{alahmadi2022falsepositives}, while Piet et al. evaluate SSH
defenses at connection and deployment scales \cite{piet2023ssh}. We performed
a dedicated search of work on hidden-state poisoning, adaptive detector
evaluation, and session- or connection-level availability. We found no prior
study that jointly measures hidden-state-poisoning detection, attacked-task
utility, adaptive evasion, and benign multi-turn session survival. This is a
scope observation used to motivate our evaluation matrix; we make no general
priority claim because adjacent security work already studies operational false
positives and connection-level effects.

\section{Spectral Foundations and Theory}

\paragraph{Condition-number convention.}
We use $\kappa=\operatorname{cond}_2(V)$ for an eigenvector matrix $V$ that
diagonalizes the state-transition operator; this is distinct from
$\operatorname{cond}_2(A)$.  In the pinned Mamba-130M checkpoint, all 24
continuous-time $A$ operators (and consequently $\exp(\Delta_t A)$) are real
and diagonal. Thus $V=I$ and $\kappa=1$ exactly. A direct audit of revision
prefix \texttt{1e76775f628f} confirms this structure; the artifact records the
complete immutable revision.
The legacy value $\kappa\approx3$ is unsupported and is not used in the
revised analysis.

Let the frozen transition at one token be diagonal:
\begin{equation}
 \bar A=\operatorname{diag}(a_1,\ldots,a_d),\qquad
 0<\rho=\max_i|a_i|<1.
\end{equation}
For $h_0\neq0$ and $0<\epsilon<1$, we define an $\epsilon$-forgetting time
using only the unforced contribution, avoiding an unobservable ``noise''
denominator:
\begin{equation}
 H_\epsilon(h_0)=\min\{t\geq0:\|\bar A^t h_0\|_2
 \leq\epsilon\|h_0\|_2\}.
\end{equation}

\begin{theorem}[Diagonal forgetting certificate]
\label{thm:diagonal-forgetting}
For $h_t=\bar A^t h_0$,
\begin{equation}
 \|h_t\|_2^2=\sum_i |a_i|^{2t}|h_{0,i}|^2
 \leq \rho^{2t}\|h_0\|_2^2,
\end{equation}
and consequently
\begin{equation}
 H_\epsilon(h_0)\leq
 \left\lceil\frac{\log(1/\epsilon)}{\log(1/\rho)}\right\rceil.
\label{eq:forgetting-upper}
\end{equation}
If $P_\rho$ projects onto coordinates with $|a_i|=\rho$ and
\begin{equation}
 \alpha=\frac{\|P_\rho h_0\|_2}{\|h_0\|_2}>0,
 \quad \|h_t\|_2\geq\alpha\rho^t\|h_0\|_2,
\end{equation}
then persistence is certified along that slow-mode component.  Spectral radius
alone never guarantees that arbitrary context persists.
\end{theorem}

The proof is an equality coordinate by coordinate and appears in
Appendix~\ref{app:diagonal-proof}.  This formulation repairs the legacy
argument, which inserted an upper bound on $\|\bar A^t h_0\|_2$ into a
signal-to-noise ratio as though it were a lower bound.  Equation
\eqref{eq:forgetting-upper} certifies forgetting, not memory capacity.

\begin{corollary}[Near-critical regime only]
\label{cor:near-critical}
If $\rho=1-\eta$ and $\eta\to0^+$, then
$-\log(1-\eta)=\eta+O(\eta^2)$, so the upper bound in
\eqref{eq:forgetting-upper} scales as $O(1/\eta)$.  This approximation is
not used for large collapse.  In particular, for $\rho=0.32$ we use the
exact denominator $\log(1/0.32)\approx1.1394$, not $\eta=0.68$.
\end{corollary}

\paragraph{Driven and time-varying systems.}
For frozen $\bar A$ and an input sequence with total energy
$\sum_k\|x_k\|_2^2\leq E$, the driven term is bounded by the finite-horizon
controllability operator; replacing it by
$\lambda_{\max}(\mathcal W_c)$ requires this energy assumption and is not
valid merely from a per-token bound.  For the selective diagonal system
$h_t=\bar A_t h_{t-1}+\bar B_t x_t$, the unforced term is exactly
$\prod_{s=1}^t\bar A_s h_0$ and satisfies
\begin{equation}
 \left\|\prod_{s=1}^t\bar A_s h_0\right\|_2
 \leq\left(\prod_{s=1}^t\rho(\bar A_s)\right)\|h_0\|_2.
\end{equation}
Every calculation using one frozen $\rho$ is therefore explicitly a local
stationary approximation, not an identity for the selective LTV system.

\paragraph{Reachability is not proved by continuity.}
Continuity of $\Delta\mapsto\rho(\exp(\Delta A))$ only describes spectral
values after a reachable $\Delta$ has been specified.  It does not imply that
the discrete token-to-$\Delta$ map is injective, surjective, or reaches an
arbitrary target.  We therefore retract the legacy ``Evasion Existence
Theorem'' and retain its optimization construction only as a falsifiable
plausibility hypothesis to be tested experimentally.

\begin{proposition}[Conditional monitor coverage]
\label{prop:conditional-coverage}
Suppose an attack crosses a declared threshold,
$\rho(\bar A_t)\leq\rho_{\min}$, at some observed token.  A monitor that
computes the diagonal maximum exactly and flags that event detects the attack
at that token.  Conversely, zero false positives are guaranteed only for a
benign class with an explicit margin
$\rho(\bar A_t)>\rho_{\min}$ at every observed token.
\end{proposition}
This proposition is conditional by construction; it makes no claim that all
effective attacks must cross the threshold and no empirical FPR claim before
the Phase~3 evaluation.

\paragraph{Exact diagonal Lipschitz constant.}
Write the continuous generator as
$A=\operatorname{diag}(-q_1,\ldots,-q_d)$ with $q_i>0$.  On
$\Delta\in[\Delta_{\min},\Delta_{\max}]\subseteq[0,\infty)$,
let $q_{\min}=\min_i q_i$.  Then
\begin{align}
 \rho(e^{\Delta A}) &= e^{-q_{\min}\Delta},\\
 |\rho(e^{\Delta_1 A})-\rho(e^{\Delta_2 A})|
 &\leq L_{\mathrm{diag}}|\Delta_1-\Delta_2|,\\
 L_{\mathrm{diag}} &= q_{\min}e^{-q_{\min}\Delta_{\min}}.
\end{align}
The coarse legacy constant
$\|A\|_2\exp(\Delta_{\max}\|A\|_2)$ ignores the negative diagonal
structure and is numerically vacuous for the measured checkpoint.  We report
$\|A\|_2=\max_iq_i$ per layer, but do not call the resulting inverse bound a
robustness certificate without a measured range for $\Delta_t$.

\section{Threat Model and Attack Construction}

The Phase~2 structural test targets the boundary of
Proposition~\ref{prop:conditional-coverage}: optimize task loss while
penalizing any decrease below a benign spectral margin.  For diagonal Mamba,
changing $\Delta_t$ cannot rotate eigenvectors---they remain coordinate axes.
The meaningful directional surface is instead the input-dependent
reachability/readout geometry through $\bar B_t$ and $C_t$.  We therefore call
the planned test a \emph{directional reachability attack}, with objective
\begin{equation}
 \max_x\;\mathcal L_{\mathrm{task}}(x)
 -\lambda\sum_t[\rho_{\min}+m-\rho(\bar A_t(x))]_+^2,
\end{equation}
and require $N\geq200$ held-out examples when token-level reachability is
tested in the experimental campaign.

\paragraph{Structural counterexample.}
We instantiated 1,000 seeded diagonal systems.  In every witness, a unit-norm
readout $C$ was rotated to be orthogonal to the retained state while $\bar A$
was held bitwise unchanged.  The remembered scalar output fell by 100\%, but a
radius-only monitor detected 0/1,000 cases (95\% Wilson upper bound 0.38\%).
For a matched spectral-collapse control, the monitor detected 1,000/1,000
cases (95\% Wilson lower bound 99.62\%).  This is a model-level counterexample
to unconditional completeness, not evidence that an arbitrary readout
rotation is reachable from discrete tokens in the pinned checkpoint.  The
latter remains an empirical question for the attack campaign.

\section{SpectralGuard Defense}

For every observed token and Mamba layer, SpectralGuard extracts the maximum
magnitude of the diagonal discretized transition and summarizes the resulting
layer-wise trajectory. The frozen detector uses 64 summary features and a
deterministic logistic regression fit on calibration data. Validation data
selects its decision threshold; test labels are never used for feature,
classifier, or threshold selection. This detector tests whether spectral
dynamics contain useful signal. It is not presented as a deployable
defense or as a complete detector.

We evaluate two response policies. Algorithm~2 is fail closed: one alert
blocks a session. Algorithm~2b is a graduated $k$-of-$w$ rule selected on
validation sessions. Both use the same frozen detector. This separation lets
us distinguish representation quality, threshold calibration, and operational
policy. It also prevents an availability improvement from being obtained by
refitting the detector on held-out sessions.

\section{Experimental Methodology}

\paragraph{Data and splits.}
The protocol covers associative recall, GSM8K, MMLU, and LongBench Qasper.
Prompt identifiers, task labels, and split assignments are frozen before
fitting. The common detector comparison uses 1,200 calibration, 800
validation, and 4,000 test rows: 1,000 held-out prompt clusters under clean,
HiSPA-original, gradient-guided, and adaptive-v2 conditions. No prompt cluster
crosses split roles. The prescribed seeds are 42, 123, 456, 789, and 2026.
The deterministic SpectralGuard fit is identical across those seed labels;
the stochastic XGBoost comparator is refit at every seed.

\paragraph{Attacks and controls.}
We retain the published HiSPA trigger as a fixed reference and add
gradient-guided, detector-aware adaptive-v2, query-only black-box,
transfer-based, and joint white-box attacks. Query and optimization budgets are
frozen before the test evaluation. A GPT-2 negative control tests for prompt
length and likelihood artifacts. Phase~5 evaluates a separately constructed
adversarial-benign trigger and a matched benign control without refitting the
detector.

\paragraph{Compute and model provenance.}
Every model load uses an immutable revision SHA, and each run records its seed,
configuration hash, package versions, model identifier, and model revision.
The confirmatory campaign evaluated 19,300 rows across seven models and four
tasks on NVIDIA H100 GPUs. Phase~3 records 26,725.9 recorded GPU-seconds across
resumable sessions. We do not add incomplete fragments from Phases~4--5 and
therefore do not report an unsupported total-compute figure. Detector fitting
uses calibration data, threshold selection uses validation data, and the test
split is evaluated once. Accuracy intervals are Wilson intervals; paired
accuracy drops and effect sizes use 2,000 bootstrap resamples. We control the
three primary paired tests with Benjamini--Hochberg at FDR 0.05 and also report
paired Wilcoxon tests.

\paragraph{Nonexchangeable uncertainty.}
For SpectralGuard, we use a 2,000-resample cluster bootstrap over the fixed test
prompt IDs, carrying all four conditions for a sampled prompt together. This
is sampling uncertainty on the fixed test set. For the CLASP-style comparator,
we report a descriptive $t$ interval over five stochastic refits: it is
seed-to-seed training variation on the same shared split, not independent
sampling uncertainty. The two interval types are labeled separately and are
not interpreted as interchangeable error bars.

\section{Results and Analysis}

The repaired GPT-2 control no longer exposes the legacy template artifact: a
classifier using only input length and prompt negative log-likelihood obtains
AUC 0.533, below the preregistered 0.65 ceiling. This result removes the basis
for the earlier categorical SSM-versus-Transformer claim; architecture-level
conclusions are deferred to the complete analysis.

Under one calibration/validation/test protocol, the multilayer SpectralGuard
detector obtains F1 0.806 and AUC 0.887 (Table~\ref{tab:detectors}). The same
fixed data and deterministic logistic fit produce identical values across the
five seed labels. Its intervals therefore come from the cluster bootstrap over
the fixed test prompt IDs, not from those duplicated fits. A CLASP-style
BOE+XGBoost reimplementation using the published design and hyperparameters,
but not official pretrained weights, obtains mean F1 0.969 and AUC 0.997. Its
interval describes seed-to-seed training variation on the same shared split,
not independent sampling uncertainty. We make no detector-superiority claim.

\begin{table}[t]
\caption{Common-protocol detector results. SG intervals are 95\% cluster
bootstrap CIs over fixed test prompt IDs; comparator intervals are descriptive
95\% $t$ intervals over five training seeds. They quantify different sources
of variation.}
\label{tab:detectors}
\centering
\small
\begin{tabular}{lcc}
\toprule
Detector & F1 & AUC \\
\midrule
SpectralGuard & .806 [.795, .818] & .887 [.878, .897] \\
CLASP-style & .969 [.966, .972] & .997 [.9962, .9968] \\
\bottomrule
\end{tabular}
\end{table}

Adaptive v2 and detector-aware v3 remain statistically indistinguishable on
the same 300 prompts: their accuracies are 1.33\% and 1.67\%, respectively
(paired difference $-0.33$ percentage points, 95\% bootstrap CI
$[-2.33,1.67]$; McNemar $p=1.0$). Logistic fits of correctness on standardized
mean spectral radius yield McFadden pseudo-$R^2$ values from 0.0089 to 0.0337,
indicating weak explanatory power for radius alone. The legacy Mamba-2.8B sign
reversal does not recur: clean-minus-perturbed $\Delta\rho$ is positive for
benign, HiSPA-original, and gradient-guided comparisons. Cross-model transfer
is reported with $N=200$ and Wilson intervals for every model and condition.

\begin{table}[t]
\caption{Adaptive task accuracy on the same 300 prompts. The comparison is
in a floor-saturated regime; failure to distinguish the attacks is not evidence of equal
power.}
\label{tab:adaptive}
\centering
\small
\begin{tabular}{lcc}
\toprule
Attack & Correct & Accuracy \\
\midrule
Adaptive v2 & 4/300 & 1.33\% \\
Detector-aware v3 & 5/300 & 1.67\% \\
\bottomrule
\end{tabular}
\end{table}

The Phase~4 escalation confirms that the conclusion is not an artifact of one
adaptive recipe. Under the frozen detector, black-box transfer, query-only,
and joint white-box attacks yield F1 0.810, 0.786, and 0.746, respectively;
their AUCs are 0.905, 0.863, and 0.829. The joint white-box cell is the
prespecified worst-case endpoint. These values show degradation under stronger
attacker knowledge, while remaining detector metrics rather than guarantees of
preserved task utility.

\section{Adaptive Threat Model}

We consider four attacker knowledge levels. The reference attacker knows the
model family and published trigger but not the detector. The gradient-guided
attacker optimizes task loss through the victim model. The query-only attacker
observes bounded detector decisions without gradients, and the joint
white-box attacker differentiates through both task and detector objectives.
Transfer attacks are generated on a surrogate and evaluated without target
model updates. All cells preserve fixed prompt sets, query budgets, attack
restart counts, and selection rules. Candidate selection occurs on validation
data; test results are not used to tune attack hyperparameters.

This protocol follows the principle that robustness claims must survive an
adaptive adversary rather than only a fixed attack \cite{athalye2018obfuscated}.
It still does not cover arbitrary token budgets, unrestricted finetuning,
weight-space attacks, or service-specific side channels. Consequently, the
joint white-box result is a strong bounded stress test, not a universal
worst-case certificate.

\section{Availability and Denial-of-Service Analysis}

We evaluated availability at the session level rather than treating a
per-window false-positive rate as sufficient.  The frozen single-alert policy
(Algorithm~2) hard-blocked all 320 clean three-turn test sessions.  A
validation-only search selected the graduated $k$-of-$w$ policy
($k=4,w=4$), which retained a 100\% hard-block rate on the frozen malicious
white-box condition.  On held-out clean sessions, however, survival increased
only to 8/320 (2.5\%, bootstrap 95\% CI [0.94\%, 4.38\%]).  Matched benign
controls and adversarial-benign prompts both yielded 0/320 surviving sessions.
The clean blocked-window rate fell from 98.74\% to 73.24\%, but monitoring
increased measured runtime by 328.53\% (4.29$\times$ total latency) over 60
validation sessions.  Thus the graduated response is directionally better than
the fail-closed baseline while remaining operationally unusable under this
frozen detector.  We report this as a negative availability result and make no
claim that Algorithm~2b resolves defense-induced denial of service.

\begin{table}[t]
\caption{Held-out three-turn session availability ($N=320$ per condition).}
\label{tab:availability}
\centering
\small
\begin{tabular}{lcc}
\toprule
Policy & Clean survival & Malicious block \\
\midrule
Fail closed (Alg.~2) & 0/320 (0\%) & 320/320 (100\%) \\
Graduated (Alg.~2b) & 8/320 (2.5\%) & 320/320 (100\%) \\
\bottomrule
\end{tabular}
\end{table}

\section{Limitations}

The forgetting theorem is exact only for a frozen diagonal transition.  The
product bound covers time variation but can be loose.  Conditional monitor
coverage does not establish that spectral collapse is necessary for task
failure.  The structural counterexample proves that a radius-only monitor is
incomplete outside its stated assumption, while not proving token-level
reachability in the learned model.
The confirmatory detector result is also weaker than the CLASP-style
reimplementation under the common protocol. The latter is a reimplementation,
not an evaluation of official pretrained CLASP weights. Low task accuracies on
several exact-answer cells limit behavioral interpretation, and the five
SpectralGuard seed labels do not create data-split independence because the
underlying split and deterministic estimator are fixed.
The availability sessions are deterministic compositions of held-out benchmark
prompts rather than live user traffic.  More importantly, the frozen detector
is nearly saturated on clean session windows: graduation improves survival only
marginally and does not distinguish the crafted benign condition from its
matched control.  Deployment therefore requires recalibration on a disjoint
session distribution or a different decision statistic; neither is evaluated
post hoc here.

\section{Conclusion}
Spectral monitoring is informative, but detection is not defense. For frozen
diagonal transitions, spectral radius certifies an upper bound on forgetting;
it does not certify arbitrary memory persistence or readout integrity. In our
experiments, richer layer-wise spectral features detect many poisoned states,
yet a CLASP-style embedding comparator is stronger on the shared split,
detector-aware attacks preserve catastrophic task failure, and both frozen
response policies impose unacceptable clean-session denial of service. The
appropriate conclusion is therefore methodological rather than architectural:
hidden-state-poisoning defenses should be evaluated jointly on detection,
attacked-task utility, adaptive evasion, and operational availability. A
deployable defense will require a statistic and calibration regime that improve
all four, not merely aggregate F1 or AUC.

\bibliographystyle{ACM-Reference-Format}
\bibliography{references}

\appendix
\section{Diagonal Forgetting Proof}
\label{app:diagonal-proof}
Because $\bar A$ is diagonal,
$\bar A^t h_0=(a_1^th_{0,1},\ldots,a_d^th_{0,d})$, hence
\begin{equation}
 \|\bar A^t h_0\|_2^2
 =\sum_i|a_i|^{2t}|h_{0,i}|^2
 \leq\rho^{2t}\|h_0\|_2^2.
\end{equation}
The right-hand side is at most $\epsilon^2\|h_0\|_2^2$ whenever
\begin{equation}
 t\geq\frac{\log(1/\epsilon)}{\log(1/\rho)},
\end{equation}
proving the upper bound after taking the ceiling.  On the dominant coordinate
set, $\|\bar A^tP_\rho h_0\|_2=\rho^t\|P_\rho h_0\|_2$; orthogonality of diagonal
coordinates gives the stated lower bound.  No controllability or noise
assumption is used.

\section{Legacy Horizon Arithmetic Audit}
The legacy pair $H(0.98)=1.17\times10^6$ and $H(0.32)=45$ cannot arise from
one formula $H(\rho)=C/\log(1/\rho)$ with fixed constants.  The implied
numerators are approximately $2.36\times10^4$ and $51.3$, respectively, a
factor above $460$.  Holding the first numerator fixed predicts about
$2.07\times10^4$ at $\rho=0.32$, not 45.  We therefore remove both values as
theoretical predictions; any future empirical horizon will be labeled and
measured directly.  The machine-readable arithmetic and layerwise norms are
generated by the Phase~2 theory-audit script.

\section{Open Science}
The linked \href{https://anonymous.4open.science/r/spectralguardaisec26/}{anonymous
review artifact} contains the
frozen configurations, split manifests, model revisions, seeds, experiment
runners, statistical analysis, tests, compact result tables, and clean
notebooks required to reproduce the paper's core claims. It records the exact
feature schema and the distinction between cluster-bootstrap intervals and
seed-to-seed descriptive intervals. Model weights and public benchmark data
are retrieved from their original providers at pinned identifiers; caches are
excluded. Credentials, machine-local state, raw multi-gigabyte downloads, and
author-identifying material are not included. Every published table is derived
from tracked compact evidence, and artifact checks reject absolute local paths,
dirty notebook outputs, unfrozen model revisions, and duplicate split roles.

\section{Ethical Considerations}
\paragraph{Scope, stakeholders, and oversight.}
This study evaluates publicly downloadable model checkpoints, public benchmark
questions, and synthetic multi-turn sessions on researcher-controlled
hardware.  It does not recruit human participants, interact with users,
inspect private or personal data, send traffic to live services, or test
operator-controlled deployments.  It is therefore not human-subjects research
and did not require institutional review board or ethics review board
(IRB/ERB) approval.  That determination does not replace substantive ethical
analysis.  The stakeholders considered are users of state-space language
models, model and defense developers, researchers who may reproduce or extend
the attacks, model providers, and the broader public affected by unreliable or
unavailable language-model services.  The analysis follows the Menlo
principles of beneficence, respect for persons, justice, and respect for law
and public interest \cite{bailey2012menlo}.

\paragraph{Risks, safeguards, and publication decision.}
The principal dual-use risk is that adaptive trigger-generation and
detector-evasion procedures could lower the effort required to disrupt an SSM
deployment.  A second risk, demonstrated directly by the experiments, is that
an over-sensitive defense can itself deny service to benign users.  The study
limits those harms by operating only on public checkpoints and offline
benchmarks; freezing bounded query and optimization budgets; excluding
credentials, service endpoints, malware, persistence, and deployment-specific
exploit automation; reporting clean-input coverage and availability alongside
attack success; and publishing negative results and scope limitations rather
than implying production readiness.  These measures do not eliminate misuse:
the attack logic remains transferable in principle.  Publication is warranted
because defenders and affected model developers gain reproducible evidence
that high token-level detection scores do not establish safe system-level
availability, while the released attacks add bounded evaluations to a
vulnerability class that is already public.  On balance, that diagnostic and
defensive benefit outweighs the residual incremental risk under the stated
safeguards.

\subsection{Responsible Disclosure and Release}
HiSPA already publicly disclosed hidden-state poisoning against Mamba-family
models \cite{lemercier2026hispa}.  This work did not probe a live product or
identify a private, operator-specific vulnerability, so there was no affected
service operator to notify before conducting the study.  The submission does
not claim a coordinated or delayed disclosure that did not occur.

The anonymous review artifact releases the defense, evaluation pipeline,
frozen configurations, compact evidence, and the bounded attack
implementations needed to verify the paper under the repository's MIT license.
It omits model-weight caches, credentials, machine-local state, and any
production targeting.  Because the attack code needed for reproducibility is
already present in the anonymous artifact, it is not described as embargoed.
If future work uncovers a previously non-public vulnerability in a live
service, the intended policy is to notify the responsible operator privately,
allow a reasonable mitigation window, and only then release details whose
defensive value justifies the remaining risk.  After review, the artifact will
be preserved at a stable archival location; acceptance would trigger
de-anonymization and replacement of the review URL with that permanent record.

\section*{Generative AI Usage Statement}
Generative AI tools were used for methodological critique, code assistance,
statistical cross-checking, literature-query formulation, and language editing.
The authors inspected and revised all generated material, executed the
experiments, verified every reported value against the released evidence, and
take full responsibility for the paper's claims, code, and citations. No
generative AI system is listed as an author.

\end{document}
